\begin{abcliste}
\abc
Prüfe jede Zeile an der Zeichnung: Bewegt sich die Springerin? Wie hoch ist sie über (3)? Ist das Seil gedehnt?
\begin{itemize}
 \item (1): in Ruhe, $\ell + s$ über (3), Seil schlaff: nur $m\,g\,(\ell+s)$. Richtig.
 \item (2): in Bewegung, $s$ über (3), Seil noch schlaff: $m\,g\,s$ und $\frac{1}{2}\,m\,v^2 = m\,g\,\ell$. Richtig.
 \item (3): Der tiefste Punkt ist ein Umkehrpunkt. Dort ist die Springerin einen Augenblick in Ruhe, $v = 0$.
\end{itemize}
\result{\text{Falsch ist } E_\mathrm{kin} = \EkbP{3}{3} \text{ in (3)}}: Richtig wäre $E_\mathrm{kin} = 0$. Die kinetische Energie von (2) ist bis (3) ganz in Spannenergie des Seils übergegangen.

\abc
Ohne Reibung ist die Summe der Energien in jeder Zeile gleich. In (1) und (2) ergibt sie \EaP{3}{3}, in (3) nach der Tabelle aber
\begin{align}
 E_\mathrm{el} + E_\mathrm{kin} &= \Ec + \Ekb = \Ef \approx \EfP{3}{3}
\end{align}
Das ist mehr, als am Anfang da war. Mit $E_\mathrm{kin} = 0$ in (3) stimmt die Bilanz: $\frac{1}{2}\,k\,s^2 = m\,g\,(\ell+s) = \EcP{3}{3}$.

\abc
Von (1) bis (2) fällt sie frei um $\ell$: $m\,g\,\ell = \frac{1}{2}\,m\,v^2$.
\begin{align}
 v_2 &= \vbF = \sqrt{2\cdot\g\cdot\lS} = \vb \approx \result{\vbP{0}{2}}
\end{align}
Sie ist dort übrigens noch nicht am schnellsten: Knapp unter (2) zieht das Seil schwächer als die Gewichtskraft, und sie wird noch schneller.
\end{abcliste}
